% questions.tex —— 1977 年·江苏扬州地区初试卷 · 题目内容片段（无答案）
% 由 tools/generate_projects.py 依来源总账自 1977/1977_数学真题_江苏扬州地区初试_试题.md 生成，勿手改（改源 MD 后重跑）


\gkpart{江苏扬州地区初试卷}

\begin{gktip}
\textbf{来源说明}：本题面依据固定源 PDF《1977江苏扬州地区初试》第 1—2 页（原书页码 245—246）逐题转录。源页为试题汇编材料中的扬州地区初试（预选）卷，仅题面、无解答；页首标题「十、扬州地区」系汇编章节标题，按原卷逐字保留为本文 H1。原卷大题以中文数字「一、」至「十、」编号，本库转写为 \textbf{1．}至 \textbf{10．}，各题分值按原卷标注在题号后（$8+8+8+8+10+10+10+10+16+12=100$ 分）。原卷无附图，各题均无小问，无附加题/参考题；配套答案与解析为本库编者独立推导，\textbf{非官方标准答案}。
\end{gktip}

\par\medskip

\begin{question}
\gkti{1}\textbf{（8 分）} 计算：

\nopagebreak[4]
\[
8^{\frac{2}{3}}+\left(\frac{1}{2}\right)^{-2}-\log_2 1+\log_3 3-\left(\sqrt{2}-1\right)^{0}.
\]

\end{question}
\begin{question}
\gkti{2}\textbf{（8 分）} 求直线 $\sqrt{3}\,x+y-2=0$ 的斜率、倾斜角及在 $x$ 轴、$y$ 轴上的截距。

\end{question}
\begin{question}
\gkti{3}\textbf{（8 分）} 把 $A$ 点的极坐标 $\left(4,\ \dfrac{3\pi}{4}\right)$ 化成直角坐标；把 $B$ 点的直角坐标 $(-\sqrt{3},\ 1)$ 化成极坐标。

\end{question}
\begin{question}
\gkti{4}\textbf{（8 分）} 不画图象求出函数 $y=2x^{2}-4x-1$ 的图象的顶点坐标、对称轴方程，并说明 $x$ 取何值时，函数取得极大值或极小值？这个极值是多少？

\end{question}
\begin{question}
\gkti{5}\textbf{（10 分）} 已知 $x$ 是正的纯小数，化简

\nopagebreak[4]
\[
\sqrt{x^{2}-2+\frac{1}{x^{2}}}\cdot\left(x+\frac{1}{x}\right)\div\frac{2x^{-2}}{1-\dfrac{x^{2}}{1+x^{2}}}.
\]

\end{question}
\begin{question}
\gkti{6}\textbf{（10 分）} 在 $\triangle ABC$ 中，$\angle A=60^\circ$，$\angle B=45^\circ$，$BC$ 边上的高 $AD=10$，求 $\triangle ABC$ 的面积（不查表，得数可保留根式）。

\end{question}
\begin{question}
\gkti{7}\textbf{（10 分）} 圆台上、下底面半径分别是 $2\ \mathrm{cm}$ 和 $6\ \mathrm{cm}$，它的侧面积是 $40\pi\ \mathrm{cm^{2}}$，求圆台体积。

\nopagebreak[4]
\[
\left[\;\left(V_{\text{\text{体}}}=\frac{1}{3}\pi h\left(R^{2}+r^{2}+Rr\right)\right)\;\right]
\]

\end{question}
\begin{question}
\gkti{8}\textbf{（10 分）} 某工厂 1977 年一月份的产值是 10 万元，如果以后平均每月产值比前一个月增长 $10\%$，问到几月份就能完成今年总产值 135 万元的任务？（$\lg 1.1=0.0414$，$\lg 2.35=0.3711$）

\end{question}
\begin{question}
\gkti{9}\textbf{（16 分）} 已知一元二次方程 $2x^{2}+4(2e-1)x+4e^{2}-1=0$ 有两个相等的实数根，以 $e$ 作为椭圆的离心率，这个椭圆的中心在原点，焦点在 $x$ 轴上，半长轴的长是 10，求这个椭圆的方程。

\end{question}
\begin{question}
\gkti{10}\textbf{（12 分）} 已知 $\triangle ABC$ 的三边 $a$、$b$、$c$ 成等差数列，$a>b>c$，$b$ 边的对角是 $\angle B$。求证：$\angle B$ 的度数小于 $60^\circ$。

\end{question}